% QMB_n11_hardware.tex — Chapter 11: Hardware
% Substantive basis: Doc. 083, 161, 183, 184; own presentation
\chapter{Hardware: Coherent Ising Machine, Google Willow, p-Bit, Photon Chip}
\label{ch:hardware}

\section{Coherent Ising Machine: Optimization as Physical Relaxation}

The Coherent Ising Machine (CIM) is an optical analog computer that
solves combinatorial optimization problems through physical relaxation
into the ground state of an Ising Hamiltonian. The fundamental principle:
The system \emph{falls} into its solution, it does not calculate it.

\subsection*{The Ising Hamiltonian as Universal Optimization Language}

Every NP-hard combinatorial problem can be formulated as minimization of the
Ising Hamiltonian:
\begin{equation}
  H_{\mathrm{Ising}} = -\tfrac12\sum_{i\ne j}J_{ij}\sigma_i\sigma_j
  - \sum_i h_i\sigma_i,
\end{equation}
where $J_{ij}$ are the coupling strengths and $h_i$ external fields.
The ground state --- the spin configuration with minimal $H_{\mathrm{Ising}}$
--- corresponds to the optimal solution. The mapping is exact, no
approximation: Travelling Salesman, MaxCut, portfolio optimization,
protein structure prediction, RSA factorization are all formulable as
Ising minimization.

\subsection*{Physical Implementation}

The CIM implements the Ising model through a network of time-multiplexed
degenerate optical parametric oscillators (DOPO):
Spin $+1$ is a laser pulse with phase $0°$, spin $-1$ a pulse with
phase $180°$; the couplings $J_{ij}$ are realized by electronic
feedback. The system evolves deterministically according to
the field equations of the parametric oscillator and thereby finds with
high probability the ground state. Up to 100,000 optical
spins are realized (as of 2024).\status{E}

A CIM appears to search stochastically; different runs sometimes yield
different solutions. But physically the system follows
\emph{deterministically} the field equations. The apparent randomness
is exclusively ignorance about the exact initial conditions of the
quantum vacuum --- a direct parallel to the FFGFT measurement.

\subsection*{Six Structural Parallels to the FFGFT}

Six parallels can be identified between CIM and FFGFT:

\textbf{Parallel 1: Determinism behind apparent stochasticity.} In
the CIM, the solution is physically enforced; in the FFGFT, the
measurement result is geometrically determined by the field configuration. Both
replace fundamental probabilism with hidden determinism.

\textbf{Parallel 2: Ground state as geometric minimum.} The CIM
relaxes into the energy minimum of the Ising Hamiltonian; the FFGFT
describes the universe as in the ground state of its own geometry
with $\xi = 4/30000$. Optimization is relaxation, not iteration.

\textbf{Parallel 3: Qubit structure.} CIM spin $\pm1$ is the phase
of the laser pulse; FFGFT qubit $|\tilde T_\uparrow,m_\downarrow\rangle$
and $|\tilde T_\downarrow,m_\uparrow\rangle$ is intrinsically enforced by $\tilde T\cdot m = 1$.
The FFGFT spin does not need to be externally limited to $\pm1$;
it is, through the geometry.

\textbf{Parallel 4: The coupling matrix.} In the classical CIM, the
$J_{ij}$ are externally programmed; in an FFGFT-based Ising machine,
they would emerge geometrically from the time field:
$J_{ij}^{(\mathrm{FFGFT})} = \int\tilde T_i\cdot\tilde T_j\cdot G(x_i,x_j)\,d^3x$.
This means: no encoding errors possible, since the couplings would be real
physical interactions.\status{S}

\textbf{Parallel 5: Fractal correction as annealing.} The correction factor
$K_{\mathrm{frak}}^{3/2}$ plays in the FFGFT the analogous role of the
annealing temperature parameter: It prevents the system from remaining in a
local minimum of the geometric energy.

\textbf{Parallel 6: Toroidal topology.} The FFGFT describes the
carrier as $T^4/\mathbb{Z}_3$. Toroidal boundary conditions would have
an advantage for a physical Ising machine: no boundary problem, all
spins have the same number of neighbors, natural implementation of
translation-invariant couplings.

\section{Google Willow: Topological Error Correction and Self-Measurement}

Google's Willow quantum processor achieved three breakthroughs in 2024 and 2025
that have special significance from the FFGFT perspective.

\subsection*{Error Correction Below Threshold (December 2024)}

Willow (105 superconducting transmon qubits, $\sim10$\,mK) shows for the first time
that the logical error rate decreases exponentially when the code distance
$d$ is increased.\status{E} A $d = 7$ block achieves an error rate by
factor $2{,}14$ below the $d = 5$ block. Physical single-qubit
error rates are $0{,}05\,\%$ --- half as much as on the
predecessor Sycamore. The error correction is based on a surface code
with topologically protected logical qubits.

FFGFT classification: More physical qubits reduce the logical error
exponentially --- this is the experimental evidence that topologically stable
configurations exist and are reachable in the quantum system, not
through external noise suppression, but through geometric order.
This is the direction in which the FFGFT points with its mode classes of the
$D_4$ lattice (Chapter~\ref{ch:gitter}).\status{S}

\subsection*{Floquet Topological Order (September 2025)}

Will et al.\ realized on Willow a Floquet topologically ordered
phase --- a quantum phase that is forbidden in thermal equilibrium,
but arises under periodic external driving.\status{E} On up to
58 qubits arose: chiral edge modes (information propagation in a
preferred direction), anyonic excitations (e-anyons become
m-anyons and vice versa), and an interferometrically measurable
topological bulk invariant.

FFGFT classification: The non-equilibrium order carries a
$\mathbb{Z}_2$ symmetry on the two-dimensional lattice --- structurally
related to the $\mathbb{Z}_3$ triality of the FFGFT carrier, but not
identical (dimension, symmetry group, and carrier geometry different).
The Floquet topological order is an experimental analogue of the
mode structure on $T^4/\mathbb{Z}_3$.\status{S}

\subsection*{Quantum Echoes (October 2025)}

Out-of-time-order correlators (OTOC) on Willow show harmonic
propagation of quantum information through constructive interference --- an
echo pattern that returns periodically. FFGFT classification: In the
torus formalism, modes are defined on a closed geometry;
information propagation returns to the source --- the echo
is a consequence of compactness, not a coincidence.

\section{The p-Bit: Between Bit and Qubit}

The p-bit (probabilistic bit) occupies the intermediate range between
the classical bit (deep potential well, $U\gg k_BT$, stable over
years) and the qubit (metastable superposition, $U\ll k_BT$, decays in
picoseconds to milliseconds). It is a magnetic nanostructure on
the geometric scale $L_8\approx22$\,nm, at which the torsion barrier
becomes thermally surmountable.

\subsection*{Physical Basis}

The energy barrier is the Stoner-Wohlfarth energy $U = K_u\cdot V$,
with the magnetic anisotropy energy $K_u$ and the volume $V$. The
switching behavior follows the Néel-Arrhenius law:
$\tau = \tau_0\exp(U/k_BT)$. At 22\,nm and room temperature, $\tau$
is in the range of microseconds to milliseconds --- not a stable bit,
but also not a coherent qubit.

The FFGFT classification: The switching behavior is deterministic on the
sub-Planck time scale $t_0 = t_P/7500\approx3{,}6\times10^{-47}$\,s;
what appears as thermal stochasticity is the macroscopically
unresolved sampling of this deterministic dynamics. The apparent
randomness is epistemic, not ontological.\status{S}

\subsection*{Reconfigurable Function}

By choice of the magnetic bias, the same p-bit can be operated as a
stochastic switch, as a forced switch, or as an analog comparator. This makes p-bits into building blocks for
probabilistic circuits (p-bit networks), which can be trained like Boltzmann
machines and tackle NP-hard problems similarly to CIMs
through physical relaxation.

FFGFT classification: The p-bit is not a quantum object (no coherent
superposition state), but a thermally fluctuating classical
mode. In the torus formalism, it corresponds to a mode whose phase $\theta$
is fully randomized by thermal noise --- the limiting case of
maximum decoherence on the carrier.

\subsection*{Why 22 nm: The Geometric Threshold}

The critical size for a p-bit at room temperature corresponds to the
FFGFT level $N = 8{,}00$, i.e., $L_8\approx22$\,nm. This is the same
threshold at which classical transistors reached the quantum limit in 2012 (Chapter~\ref{ch:spin}). Three independent phenomena converge
on the same geometric scale: transistor limit, excitonic
Bohr radius, p-bit threshold.\status{S}

\section{Photon Chip: Room-Temperature Quantum Photonics}

China's CHIPX/Touring-Quantum chip (2025) is a monolithic
6-inch thin-film lithium niobate (TFLN) wafer with over 1000 optical
components. Key data: $1000\times$ speedup for parallel tasks,
$100\times$ lower energy consumption, stable at room temperature,
$12000$ wafers per year production.\status{E}

\textbf{Why room temperature works:} Optical photons have
$E\approx1$--2\,eV; $E/k_BT(300\,\text{K})\approx40$. Thermal
noise cannot create optical photons; the system sits far
below the thermal threshold (Chapter~\ref{ch:qubit}). The TFLN
single crystal with $\sim10^{23}$ lattice sites is geometrically extremely
robust.\status{K}

\textbf{FFGFT classification:} Massless modes (photons) are in the FFGFT
three-orbits in $GF(27)^*$ --- structurally separated from the massive fermions
(fixed points $\{+1,-1\}$) (Frobenius separation, Doc.~339).\status{B}
The photonic platform is suitable for loophole-free Bell tests and
is thus the natural testbed for the $\xi$ deviation
$\Delta\mathrm{CHSH}\approx\xi/(2\pi)\approx2\times10^{-5}$ --- if
detection efficiency and coherence times suffice.

\textbf{Proposed FFGFT optimizations:} Three approaches are conceptually
sketched, all not yet realized: (1) Topology compiler that performs
topological qubit placement according to the $D_4$ lattice structure
and could reduce SWAP gates by 30--50\,\%; (2) harmonic resonance
with qubit frequencies at $\xi$ multiples; (3) time field modulation for
active coherence preservation. Status: [S] for all three.

\quelle{Doc.~083 (Photon chip 2025), Doc.~161 (Coherent Ising Machine), Doc.~183 (Google Willow, April 2026), Doc.~184 (p-Bit), Doc.~339 (Frobenius separation), Doc.~343 (Resolution floor)}